\chapter{Introduction}
% =========================================================
\label{sec:intro}
\noindent
Modern \emph{sequence generators} are the engines behind a wide range of applications: they translate sentences from one language to another in \gls{mt} systems, turn audio into text in \gls{asr} systems, and produce descriptions of images in image captioning, among many other usages. In all of these settings, a model proposes many possible outputs token by token, and a \emph{decoder} must decide which full sequence to present. The dominant practice is to prefer what the model believes is most likely. Procedures such as greedy search and beam search operationalize this idea by approximating \gls{map} decoding---selecting the single hypothesis with the highest model probability.

Over the last few years, evidence has mounted that probability alone is an imperfect stand-in for what humans actually value. In machine translation, for instance, sentences judged better by neural evaluation metrics like BLEURT can receive \emph{lower} probability than the beam-search winner \cite{mbr_neural_metrics_2021,eikema2021beyondbeam}. Similar mismatches appear in \gls{asr}, where the highest-confidence transcript is not always the one readers perceive as most accurate. These observations motivate a shift from likelihood-first decisions to \emph{quality-first} ones.%, as illustrated in Figure~\ref{fig:likelihood_vs_quality}.

\iffalse
\begin{figure}[t]
  \centering
  \includegraphics[width=0.6\linewidth]{figures/likelihood_vs_quality.png}
\caption{Motivation: likelihood-first versus quality-first selection. Given an input \(x\) (text, audio, or image), a conditional generator \(p_{\theta}(y\mid x)\) assigns probabilities to many candidate outputs \(y\). Conventional decoding (e.g., MAP approximation via greedy or beam search) selects and returns \(y_{\mathrm{MAP}}\), the single hypothesis with the highest model probability. The same candidate set can also be evaluated under a task-aligned quality signal \(u\), representing either a reference-based metric (e.g., \(\mathrm{WER}(y,r)\) in ASR or CIDEr against COCO references) or a reference-free scorer (e.g., CLIP-based image--text alignment). Because model likelihood is optimized for next-token prediction rather than for the downstream notion of quality, the probability ranking and the quality ranking can disagree: a lower-probability candidate may receive a higher utility and be preferred by humans. Quality-aware decoding exploits this gap by generating multiple plausible candidates and selecting the output that is best under \(u\) (or under an expected-utility criterion such as MBR).}
  \label{fig:likelihood_vs_quality}
\end{figure}
\fi
This work proposes the development and evaluates \gls{qad} strategies as general alternatives for different fields. The core idea is straightforward: instead of committing to a single output, we can generate several plausible candidates and then choose the one that scores highest under a task-appropriate notion of quality \cite{qad_nmt_2022}. What counts as ``quality'' depends on the application. For \gls{asr} it means fewer mistakes and better preserved meaning, typically captured by metrics such as the \gls{wer} or learned referenceless metrics. For image captioning it means relevance and semantic faithfulness to the scene, measured by lexical or neural metrics. Across modalities, the common pattern is a generate--and--re-rank pipeline driven by the utility that best reflects human judgement.

A principled way to formalize this shift is through \gls{mbr} decoding. Rather than asking which output is most probable, \gls{mbr} asks which output has the best \emph{expected} utility when compared to plausible alternatives \cite{kumar_byrne_mbr_smt_2004,tromble_lattice_mbr_2008,bertsch_mbr_2023}. In practice, this expectation is approximated by constructing an \(N\)-best list or a set of samples and then evaluating each candidate with the chosen utility. When the utility aligns with human preferences, \gls{mbr} can improve upon other types of decoding (such as MAP) \cite{qad_nmt_2022,mbr_neural_metrics_2021,bertsch_mbr_2023}.

Quality gains come with computational choices, so \emph{test-time scaling} is also examined. The number of candidates, their \emph{diversity}, and model size trade-off in terms of accuracy and cost will be assessed. Intuitively, more—and more varied—candidates will increase the chance that at least one is excellent and provide stronger evidence for consensus during re-ranking. The resulting error can thus shrink exponentially or as a power law with the number of candidates \cite{bertsch_mbr_2023}. This leads to a practical question that will be empirically addressed: assessing whether or not smaller or mid-sized models, paired with richer candidate sets and good utilities can rival the quality of larger models \cite{test_time_scaling_reranking_2025}: by exploring these trade-offs across \gls{asr} and image captioning tasks, this work reframes decoding as a quality-first decision and offers guidance on how to allocate computation where it most improves human-perceived performance.

The remainder of this document is organized as follows. Chapter~2 introduces the foundational concepts to be developed and used throughout the work, including conventional decoding strategies, candidate generation and diversity, reranking, and the decision-theoretic motivation behind \gls{mbr}. Chapter~3 reviews related work on quality-aware decoding in \gls{mt}, \gls{asr}, and image captioning. Chapter~4 presents the methodology adopted in this work, describing the experimental setup, the candidate generation strategies used to control diversity, the utilities and evaluation metrics employed for selection, and the protocol used to study quality--compute trade-offs under test-time scaling.
